\chapter{Análise de sensibilidade: vazamento de alvo e partição por REFREC}\label{ap:sensibilidade}

Os indicadores de encerramento por óbito do Sinan-TB e a presença de potenciais pares associados ao mesmo registro do SIM em treino e teste constituem limitações distintas dos Estudos 2 e 3 (Seções~\ref{sec:e2-limitacoes} e~\ref{sec:e3-limitacoes}). A análise deste apêndice separa a contribuição desses dois componentes na ablação e no impacto epidemiológico, preservando os resultados principais como registro dos experimentos originais. No \textit{linkage}, também denominado relacionamento de registros (\textit{record linkage}), a avaliação de potenciais pares não equivale à avaliação de registros independentes, o que exige explicitação da unidade de partição \cite{Jardim2026phdrepo}.

\section{Método}

Foram cruzadas duas escolhas: a partição estratificada por linha ou agrupada por REFREC, identificador do registro do SIM, e o vetor original de 58 atributos ou o vetor de 56 atributos, obtido pela retirada de \texttt{obito\_tb} e \texttt{obito\_sinan}. Definiram-se os cenários A (linha, 58), B (linha, 56), C (REFREC, 58) e D (REFREC, 56). A matriz reduzida foi derivada da matriz canônica, sem alteração da engenharia dos demais atributos. Os cenários A e B compartilham as mesmas linhas de treino e teste; C e D compartilham os mesmos grupos, sem REFREC em comum entre treino e teste. A fração nominal de teste foi mantida em 0,3 e a semente em 42; na partição agrupada, a estratificação considera a presença de ao menos um vínculo no grupo, não a proporção de vínculos em cada linha.

Em cada cenário, treinaram-se Random Forest com pesos de classe, Random Forest com SMOTE (RF+SMOTE), Gradient Boosting (GB) e XGBoost local (CPU). Foram preservados os hiperparâmetros da ablação para os três primeiros; o SMOTE, com razão de 0,3 e três vizinhos, foi aplicado somente ao treinamento do RF+SMOTE. O XGBoost local utiliza os hiperparâmetros do fluxo Modal, com execução em CPU, sem presumir identidade numérica com o resultado produzido naquele ambiente. As 70 configurações originais, distribuídas em nove categorias de decisão, foram avaliadas pela mesma função reutilizável da ablação, sem incorporar o XGBoost local ao consenso original. Para este último, reporta-se o ponto de probabilidade de 0,5. Os melhores valores de $F_1$ por categoria permanecem comparações exploratórias no teste, não estimativas independentes de seleção de modelos.

O impacto epidemiológico foi calculado para o RF+SMOTE com probabilidade de pelo menos 0,5, em comparação à nota final de pelo menos 8. Denominam-se recuperados os vínculos identificados pelo modelo e não pela nota; o incremento líquido corresponde à diferença entre o total detectado pelo modelo e pela nota, descontando os vínculos encontrados apenas por esta última. A variação percentual tem como denominador os vínculos detectados pela nota, não o total de vínculos do teste. São contagens entre potenciais pares rotulados; a mudança da partição altera a composição do conjunto avaliado, impedindo interpretar a diferença bruta entre cenários como efeito causal sobre a mortalidade.

Os agregados encontram-se em \texttt{data/sensibilidade\_sem\_obito\_refrec.csv}, e os metadados, no arquivo de mesmo nome terminado em \texttt{\_meta.json}, com versões, hiperparâmetros, resumos das configurações, tempos e resumos criptográficos das entradas e dos códigos. A execução por \texttt{scripts/sensibilidade\_sem\_obito\_refrec.py} verifica individualmente os verdadeiros positivos, falsos positivos e falsos negativos das 70 configurações do cenário A contra \texttt{data/ablation\_results.csv}. Nenhum identificador ou linha de potencial par integra essas saídas \cite{Jardim2026phdrepo}.

\section{Resultados e interpretação}

% Gerado por scripts/gerar_tabela_sensibilidade.py de data/sensibilidade_sem_obito_refrec.csv
\begin{table}[htbp]
\centering\footnotesize
\setlength{\tabcolsep}{3pt}
\caption{Sensibilidade da ablação e do impacto epidemiológico à retirada dos indicadores de óbito e à partição por REFREC.}
\label{tab:sensibilidade-2x2}
\begin{tabular}{lrrrrrrr}
\toprule
Cenário & $n$ teste & Vínculos & \shortstack{$F_1$ RF\\+SMOTE} & $F_1$ GB & \shortstack{$F_1$ XGB\\local} & \shortstack{VP/FP\\RF+SMOTE} & \shortstack{Recuperados/\\líquido} \\
\midrule
A: linha/58 & 18.509 & 74 & 0,931 & 0,870 & 0,931 & 67/3 & 27/24 \\
B: linha/56 & 18.509 & 74 & 0,938 & 0,867 & 0,931 & 68/3 & 27/25 \\
C: REFREC/58 & 18.304 & 72 & 0,846 & 0,824 & 0,869 & 63/14 & 28/25 \\
D: REFREC/56 & 18.304 & 72 & 0,846 & 0,816 & 0,875 & 63/14 & 28/25 \\
\bottomrule
\end{tabular}
\par\smallskip\begin{minipage}{\textwidth}\footnotesize
Nota: RF: Random Forest; GB: Gradient Boosting; XGB: XGBoost local (CPU), não o resultado do fluxo Modal. O limiar de probabilidade de 0,5 e a nota $\geq 8$ são os da ablação. Recuperados são os vínculos encontrados pelo RF+SMOTE e não pela nota; líquido é a diferença entre as detecções. Com agrupamento por REFREC, o denominador e os vínculos avaliados mudam. 56 atributos excluem \texttt{obito\_tb} e \texttt{obito\_sinan}. Fonte: \texttt{data/sensibilidade\_sem\_obito\_refrec.csv}.
\end{minipage}
\end{table}


Na partição por linha, os cenários A e B avaliaram 18.509 potenciais pares, com 74 vínculos e 1.184 registros REFREC distintos no teste. A retirada dos indicadores de óbito elevou o $F_1$ do RF+SMOTE de 0,931 para 0,938, com passagem de 67 para 68 verdadeiros positivos e manutenção de três falsos positivos. O total de vínculos recuperados abaixo da nota 8 permaneceu em 27, mas a interseção com a nota passou de 40 para 41; frente aos 43 vínculos detectados pela nota, o incremento líquido passou de 24 para 25, ou de 55,8\% para 58,1\%. O GB apresentou pequena redução de $F_1$, de 0,870 para 0,867, enquanto o XGBoost local manteve 0,931. Nessas condições, o efeito da retirada dos indicadores foi marginal e não houve redução do desempenho do RF+SMOTE nem do seu impacto líquido.

Com agrupamento por REFREC, os cenários C e D avaliaram 18.304 potenciais pares, com 72 vínculos e 454 registros REFREC distintos no teste. O $F_1$ do RF+SMOTE foi de 0,846 em ambos, com 63 verdadeiros positivos e 14 falsos positivos, ante 67 e três no controle A. A nota detectou 38 vínculos, e o modelo, 63, dos quais 35 foram detectados pelos dois critérios e 28 foram recuperados apenas pelo modelo; o incremento líquido foi de 25, equivalente a 65,8\% das detecções pela nota. O GB passou de $F_1$ de 0,824 em C para 0,816 em D, e o XGBoost local, de 0,869 para 0,875. A redução de $F_1$ do RF+SMOTE de 0,931 em A para 0,846 em C, e de 0,938 em B para 0,846 em D, foi mais acentuada que a variação decorrente da retirada dos indicadores. A partição por linha apresentou 1.086 registros REFREC compartilhados entre treino e teste, contra nenhum nas partições agrupadas; contudo, a mudança dos vínculos avaliados e do denominador impede atribuir toda a diferença exclusivamente à dependência entre os conjuntos.

Na importância por ganho médio normalizado do XGBoost local, \texttt{obito\_tb} e \texttt{obito\_sinan} ocuparam, respectivamente, as posições 22 e 18 em A, e 54 e 25 em C. Em C, o ganho de \texttt{obito\_tb} foi zero, e sua posição decorre do desempate pela ordem das colunas entre atributos sem ganho. Esse ordenamento não é o ranking SHAP do experimento original. Os resultados apontam para efeito limitado do vazamento declarado pelos indicadores de óbito nas configurações aqui avaliadas, sem tornar legítimo seu uso como preditores; o otimismo do resultado principal do RF+SMOTE se manifesta sobretudo no contraste com a partição agrupada, com aumento dos falsos positivos e redução de $F_1$, embora se mantenha o incremento líquido de vínculos detectados. Essa leitura restringe-se ao experimento de sensibilidade e não se estende ao QW1 nem constitui estimativa de efeito sobre a mortalidade populacional \cite{Jardim2026phdrepo}.

\section{Limites da análise}

A retirada dos indicadores de óbito e o agrupamento por REFREC não corrigem a escolha histórica dos limiares no conjunto de teste. Os melhores resultados por categoria continuam sujeitos a otimismo de seleção, e uma única partição não estima a variabilidade entre amostras. O agrupamento impede a repetição do registro do SIM, mas não garante ausência de dependência por registro do Sinan-TB, identificado por COMPREC. O XGBoost local não substitui o resultado do fluxo Modal, e a importância por ganho não equivale à atribuição SHAP. O QW1 não foi reexecutado, por depender de GPU; seus resultados e os da ablação de medidas contínuas permanecem sujeitos às ressalvas dos respectivos protocolos. A análise não modifica o GZ-CMD, que não utiliza os indicadores de óbito e já dispõe de calibração fora da amostra agrupada por REFREC, nem dispensa a validação externa em outra população \cite{Jardim2026phdrepo}.
